\chapter{Kết luận}

Dựa trên mục tiêu ban đầu và phạm vi nghiên cứu đã xác định, nhóm đã chuẩn bị kỹ lưỡng và hiện thực được các chức năng nền tảng của một hệ thống ôn luyện tiếng Anh ứng dụng AI. Về mặt chức năng, nền tảng đáp ứng được các yêu cầu cốt lõi của một ứng dụng web ôn thi tiếng Anh hiện đại, bao gồm hỗ trợ nhiều kỳ thi chuẩn hóa như IELTS và TOEIC, đồng thời cung cấp môi trường thi thử sát thực tế. Điểm đột phá quan trọng nhất của đề tài là việc thiết kế thành công một quy trình đánh giá kỹ năng Nói toàn diện. Bằng cách kết hợp các công nghệ như OpenAI Whisper cho nhận dạng giọng nói, Librosa cho phân tích độ trôi chảy, Parselmouth cho đánh giá ngữ điệu và GPT-4o-mini cho phân tích ngôn ngữ, hệ thống đã cung cấp cơ chế chấm điểm khách quan cùng phản hồi chi tiết.

Mặc dù đã đạt được nhiều kết quả tích cực, nhóm thực hiện cũng thẳng thắn nhìn nhận một số hạn chế còn tồn tại của hệ thống. Cụ thể, mô hình xử lý tuần tự trong quy trình đánh giá kỹ năng Nói tuy đảm bảo ổn định nhưng chưa được tối ưu hoàn toàn về độ trễ thời gian thực. Ngoài ra, sự phụ thuộc vào API của bên thứ ba cũng đặt ra thách thức về chi phí vận hành khi mở rộng quy mô. Dù vậy, xét trên tổng thể, dự án đã thành công trong việc xây dựng một nền tảng giáo dục tích hợp sâu với trí tuệ nhân tạo, không chỉ dừng lại ở việc số hóa bài thi mà còn tạo ra giá trị gia tăng thông qua các công cụ đánh giá tự động và tư vấn học tập, làm tiền đề vững chắc cho các hướng phát triển tiếp theo.